% 第五章正文，直接编辑本文件。
\FloatBarrier\clearpage

\section{总结}\label{chapter:5}

\subsection{工作总结}\label{section:5.1}

本作品针对信息茧房难以被个人察觉、定位和核查的问题，设计并实现了个人认知体检工具``破茧''。系统以手机和浏览器端的增量采集为入口，记录用户实际观察的内容片段与接触过程；在可信执行环境内完成多模态观点与论据联合理解、面向个人历史的信息增量识别和基于时序论证图的长期接触模式分析；结果以小时报告、日报和长期报告呈现，每项结论均可回到具体原文片段核查，用户可对错误归属和采集范围作出更正。

围绕四项研究任务，本报告分别给出了方法设计与实验验证。内容理解方面，已观察前缀建模、目标条件对齐与质量门控使分析范围与用户实际接触保持一致，论证字段 macro-F1 为 0.846，主体归属准确率为 0.934，新增片段到达后结果按版本修订。历史比较方面，观点与论据双路表示配合多历史联合覆盖，Recall@50 为 0.941，五维增量 macro-F1 为 0.837，能够区分``相同观点获得新依据''与``相同观点反复换写''。长期识别方面，时序论证图在用户隔离协议下 macro-F1 为 0.834、ECE 为 0.034，正常专题学习的误报率降至 6.4\%。安全与运行方面，远程证明、密钥分离、检查点防回退与删除重放构成完整的机密计算协议，规定的恶意用例全部被拒绝或隔离，机密执行相对常规执行的分析时延增量为 11.18\%；系统在 10 事件/秒的基础负载下长期稳定运行。

与既有方案相比，作品的差异集中在三处：分析对象限定为个人授权范围内可核查的浏览经历，而非群体结构或问卷自报；判别依据采用观点重复、五维信息增量与时间持续性的组合，而非主题占比或文本相似度；隐私保护覆盖配置一致性、派生数据与历史状态管理，而非仅保护单次推理的输入。上述差异的实际作用均由第三章的对照与消融实验分别检验。

\subsection{不足与适用范围}\label{section:5.2}

系统仍存在明确的边界。采集层在极快切换和平台页面结构变化时仍有误报与漏报，跨应用采集依赖系统能力与目标应用的授权策略；语义层对反讽、多主体同屏和隐含条件的处理不完善，未观察区域的内容不进入分析；长期结论受观察覆盖限制，采集不足的窗口只能输出``证据不足''；情绪与个体反应的一致性基于自愿自报反馈，不构成临床效度，也不证明浏览内容与个人情绪之间的因果关系。

报告的性质是用户审视自身信息环境的材料。系统不将内容相似、观点一致或负面情绪直接作为认知受损的证据，不根据浏览记录推断用户认同或心理状态；对用户信息选择的最终帮助，需要后续用户研究检验。本章所述性能与安全结论均对应第三章规定的硬件、负载与威胁模型，不外推为一般条件。

\subsection{后续研究方向}\label{section:5.3}

后续工作可从三个方向展开。一是用户研究：招募更大规模的参与者，检验认知体检报告是否改善用户对自身信息局限的觉察，以及这种觉察是否转化为信息选择行为的变化。二是方法改进：针对残余误差较集中的反讽识别、隐含条件恢复和长视频部分观察场景扩充训练数据，并探索覆盖不足时更细粒度的拒识策略。三是工程完善：扩展 iOS 与更多浏览器平台的采集适配，降低机密执行在会话建立与状态提交环节的固定开销，并研究多设备并行提交下的协调协议。
